% =======================================================
% 4. DEFINITION OF DONE
% =======================================================
\section{Definition of Done (DoD)}
\label{sec:dod}

La Definition of Done\textsubscript{G} definisce i criteri minimi che un task deve
soddisfare per essere considerato completato.

I criteri operativi della Definition of Done sono definiti nelle \link{https://synergo-unit-unipd.github.io/Second-Brain/docs/RTB/doc_gestionali/doc_interni/norme_di_progetto/norme_di_progetto.pdf}{Norme di
Progetto}, all'interno del processo di Verifica. Nel presente Piano di Progetto
la DoD viene richiamata come criterio di chiusura dei task e degli sprint,
evitando di duplicare regole operative già normate nel way of working del
gruppo.

Un task è considerato \textit{Done} solo se soddisfa:
\begin{itemize}
\item i criteri comuni previsti dalle Norme di Progetto;
\item i criteri specifici per task documentali, nel caso di attività di
documentazione;
\item i criteri specifici per task di implementazione del Proof of
Concept\textsubscript{G}, nel caso di attività di sviluppo del PoC.
\end{itemize}

La verifica del rispetto della DoD avviene durante la Pull Request\textsubscript{G}
e viene confermata in Sprint Review\textsubscript{G}. Eventuali task non
completati vengono riportati nel backlog\textsubscript{G} e motivati nel
verbale interno dello sprint.

% ── 4.1 Livello task ────────────────────────────────────────────────────────
\subsection{Livello task}
\label{sec:dod-task}

A livello di task, la DoD viene applicata in base alla natura dell'attività.

Per i task documentali, il completamento richiede che il contenuto sia coerente
con lo scope della issue\textsubscript{G}, che il documento compili correttamente,
che eventuali link e riferimenti interni siano verificati, che il Glossario sia
aggiornato se necessario e che la Pull Request abbia ricevuto almeno una review
positiva da un Verificatore\textsubscript{G} diverso dall'autore della modifica.

Per i task di implementazione del PoC, il completamento richiede che il codice
sia coerente con lo scope della issue, che il PoC si avvii localmente, che non
siano presenti errori evidenti in console, credenziali, token o file sensibili
nel repository\textsubscript{G}, e che eventuali modifiche rilevanti siano
riflesse nei documenti coinvolti.

% ── 4.2 Livello sprint ──────────────────────────────────────────────────────
\subsection{Livello sprint}
\label{sec:dod-sprint}

Uno sprint è considerato \textit{Done} se:
\begin{itemize}
\item tutti i task pianificati sono conclusi secondo la DoD, oppure sono
stati esplicitamente rimandati al backlog con motivazione riportata
nel verbale interno;
\item la Sprint Review ha verificato gli output prodotti rispetto agli
obiettivi dello sprint;
\item la Sprint Retrospective\textsubscript{G} ha individuato eventuali
criticità e, se necessario, almeno un'azione correttiva;
\item le eventuali azioni correttive individuate durante la Sprint
Retrospective sono recepite nel processo di miglioramento descritto
nelle Norme di Progetto e tracciate tramite issue o verbali interni;
\item il registro dei rischi è stato rivalutato;
\item le ore effettive sono state rendicontate;
\item eventuali decisioni o modifiche operative emerse sono state tracciate
tramite issue o verbale.
\end{itemize}
